\section{Experimental Methodology}
\label{sec:method}
% NEW (written for the TOMPECS version; not in the original paper)
Every result in this paper is a measured execution time on shared hardware, so its credibility
rests on the protocol stated here: what was run, on what, how correctness was established, and
how time was measured and summarized.

\subsection{Workloads, Models and Platforms}
\label{sec:explanation}
We run two families of experiments: a prompt-component ablation on single-kernel benchmarks, and an
end-to-end study on production applications.

\paragraph{Benchmarks.}
The prompt ablation uses six HeCBench \cite{jin2023hecbench} CUDA benchmarks---\texttt{atan2},
\texttt{bitonic-sort}, \texttt{md5hash}, \texttt{accuracy}, \texttt{sobol} and
\texttt{shmembench}---chosen to span memory-bound and compute-bound behaviour and to include kernels
with built-in verification. The profile-data ablation (Section~\ref{sec:profabl}) uses a further eleven
HeCBench applications: \texttt{adv}, \texttt{amgmk}, \texttt{atomicSystemWide},
\texttt{background-subtract}, \texttt{bincount}, \texttt{blockAccess}, \texttt{boxfilter},
\texttt{ccsd-trpdrv}, \texttt{chemv}, \texttt{concat} and \texttt{convolutionSeparable}.

\paragraph{Applications.}
The end-to-end study uses two production codes: AutoDock-GPU (molecular docking, CUDA) and Cholla
(astrophysical hydrodynamics, CUDA). Each is built
from its own build system with its own inputs, and validated by a gate derived from its scientific
output rather than from a benchmark harness (Section~\ref{sec:realapps}).

\paragraph{Models.}
The reported prompt-ablation grid uses GPT-5 and Gemini~2.5~Pro. Every later study---the
semantic-preservation constraint, the applications, the profile-data ablation and the
prescriptive-guidance comparison---uses GPT-5 and Gemini~3.1~Pro; we note the model change explicitly wherever numbers from the two
studies appear together, since they are not directly comparable.
% Earlier exploratory runs used llama3.1:70B, llama2:7B and qwen-coder:14B. They are not in any
% reported table; add a paragraph here if you want to include those numbers.

\paragraph{Hardware.}
% EDIT (2026-10-26): the prescriptive-guidance benchmarks (Section 8.2) were timed on the GH200, not
% the H100 (four_arm_hopper2.csv).
The prompt ablation, the semantic-preservation comparison and the HeCBench part of the profile-data
ablation are timed on an NVIDIA H100 PCIe (sm\_90). The application study, the real-application part
of the profile-data ablation and the prescriptive-guidance comparison of Section~\ref{sec:guidvar}
are timed on an NVIDIA GH200 (sm\_90). Both nodes use CUDA 12.6.3. One
mismatch is worth stating: the profiles supplied in the prompt for the six original ablation benchmarks
were collected on an A100 ($108$ SMs) rather than on the H100 ($114$ SMs) they were then timed on; the
eleven profile-data-ablation benchmarks use native H100 profiles. All nodes are shared with other
users, which motivates the measurement protocol of Section~\ref{sec:timing}: build once, interleave
variants within one session, wait for an idle GPU, and log competing GPU processes at the start and end
of every repetition so that contaminated repetitions can be excluded rather than averaged in.

\subsection{Workload Selection}
The benchmark study spans memory-bound (\texttt{shmembench}) and compute-bound (\texttt{bitonic-sort},
\texttt{md5hash}, \texttt{atan2}) kernels. The application study selects, for each program, the single
file containing the kernel that dominates GPU time, and---critically---only after confirming that this
kernel also dominates \emph{application} time.

\subsection{Correctness Validation}
Correctness is checked in whichever of two ways the program supports. If the program prints numerical
results, the optimized version must reproduce the baseline's numbers; if it ships a built-in
verification, we use its verdict, tailored per application. For the production applications neither is
sufficient on its own, so we construct explicit gates---bit-exact field comparison where the
application is deterministic, and a statistical band where it is not---and validate each gate by
injecting a deliberate fault and confirming the gate rejects it (Section~\ref{sec:realapps}).

\subsection{Timing and Measurement Protocol}
\label{sec:timing}
Time comes from the program's own reported kernel or phase time wherever it provides one, and from wall
clock otherwise; for the applications this matters, since AutoDock-GPU's CUDA setup alone is $68$\,s on
a cold driver and would swamp a few percent of solver time. The number reported by the pipeline during
a run is indicative only. Every number we report, except the prompt-component grid of
Section~\ref{sec:prompts} (whose runtimes are the pipeline's own ten-run averages), comes from a
separate measurement pass that builds
each arm once, discards a warmup run, interleaves all arms of an application round-robin within a
single session, waits for an idle GPU before starting, and logs other users' GPU processes at the start
\emph{and} the end of every repetition. We take medians over the repetitions that were clean by that
test---five per arm for the benchmarks, and for the applications between two and seven, depending on
how often other users' jobs intervened (each table states its count)---and treat repetitions that disagree with their neighbours as suspect even when no foreign
process was seen. Between-session drift on these shared nodes reached $49\%$ on identical code, which
is an order of magnitude larger than most of the effects we are trying to measure.

% NEW (written for the TOMPECS version; not in the original paper)
\subsection{Reporting Conventions}
Speedups are reported against the unmodified program unless a table states otherwise, and every
reported number outside Section~\ref{sec:prompts} comes from the separate measurement pass described
above rather than from the pipeline's own in-run timing. Except where stated, each (application, model, prompt) cell is a
single language-model sample, and an effect below roughly $5\%$ in a single cell is treated as
unresolved rather than small (Section~\ref{sec:limits}). Where two arms are compared on a noisy
benchmark, the comparison is scored against that pair's own repetition spread
(Section~\ref{sec:profabl}), and any speedup above $5\times$ is treated as a correctness alarm and
its diff inspected (Section~\ref{sec:intent}).
